\subsection{点测度. 有限支集的测度}

命题 12. --- 设 \(a_i\) （ \(1 \leqslant i \leqslant n\) ）是局部紧空间 X 中互不相同的点。X 上每个支集含于诸 \(a_i\) 所成之集的测度都是诸测度 \(\varepsilon_{a_i}\) （ \(1 \leqslant i \leqslant n\) ）的一个线性组合。

因为这样的测度 \(\mu\) 对每个函数 \(f\in\mathcal{K}(\mathrm{X};\mathbf{C})\) （满足 \(n\) 个关系 \(f(a_{i})=0\) ）都为零（No.~3, 命题 8）；由于这些关系可写成 \(\varepsilon_{a_{i}}(f)=0\) ，故 \(\mu\) 是诸 \(\varepsilon_{a_{i}}\) 的线性组合（A, II, §7, No.~5, Th. 7 的推论 1）。

特别地，每个支集或为空集或退缩为单独一点 x 的测度都具有 \(\alpha\varepsilon_{x}\) 的形式，其中 \(\alpha\) 是一个复数；这样的测度称为点测度；于是每个支集为有限集的测度是点测度之和。

定理 1. --- 每个测度 \(\mu\) （局部紧空间 \(X\) 上的）都属于向量空间 \(V\) （由支集有限且含于 \(\operatorname{Supp}(\mu)\) 的测度构成）的淡闭包。

只需证明 \(\mu\) 与子空间 \(V^{\circ}\) （\(\mathcal{K}(X; \mathbf{C})\) 中与 \(V\) 正交的子空间）正交（TVS, II, §6, No.~3, Th. 1 的推论 2），也就是说，关系式 \(\langle f, \varepsilon_{a} \rangle = 0\) （其中 \(a\) 跑遍 \(\mu\) 的支集）蕴涵 \(\langle f, \mu \rangle = 0\) ；而这正是 No.~3 的命题 8。

推论 1. --- X 上的每个有界测度 \(\mu\) 都属于由支集有限、含于 \(\mu\) 的支集且范数 \(\leqslant \| \mu \|\) 的测度所构成的凸集 A 的淡闭包。此外，若 \(\nu\) 保持在 A 中而淡趋于 \(\mu\) ，则 \(\| \nu \|\) 趋于 \(\| \mu \|\) 。

为证第一个断言，只需证明测度 \(\mu\) 属于极集 \(A^{\circ}\) （A 在 \(\mathcal{K}(X;C)\) 中的极集）的极集（TVS, II, §6, No.~3, Th. 1 与 §8, No.~4）；这就是说，对 \(f \in \mathcal{K}(X;C)\) ，关系 \(|\langle f,\varepsilon_{a}\rangle| \leqslant 1/\|\mu\|\) 对一切 \(a \in \operatorname{Supp}(\mu)\) 成立蕴涵 \(|\langle f,\mu\rangle| \leqslant 1\) ；而这是 No.~3 命题 8 的推论 3 的推论。

为证第二个断言，我们注意

\[
\liminf _ {\nu \to \mu ,   \nu \in A} \| \nu \| \geqslant \| \mu \|
\]

因为函数 \(\nu \mapsto \| \nu \|\) 对淡拓扑是下半连续的（§1, No.~9, 命题 15 的推论 4），而结论由如下事实得出：按定义， \(\| \nu \| \leqslant \| \mu \|\) 对 \(\nu \in A\) 成立。

推论 2. --- 每个有界测度 \(\mu\) （\(X\) 上的）都属于由支集有限、含于 \(\mu\) 的支集且范数等于 \(\|\mu\|\) 的测度所成之集的淡闭包。

可以假设 \(\mu \neq 0\) 。设 \(V\) 是淡拓扑下 0 的一个开邻域；对每个 \(\varepsilon\) （满足 \(0 < \varepsilon < 1\) ），依据推论 1，存在一个测度 \(\nu_0\) ，其支集有限且含于 \(\operatorname{Supp}(\mu)\) ，并使 \(\nu_0 - \mu \in V\) 且 \(\| \mu \| \geqslant \| \nu_0 \| \geqslant (1 - \varepsilon) \| \mu \|\) 。令 \(\nu = (\| \mu \| / \| \nu_0 \|) \nu_0\) ，则有 \(\| \nu \| = \| \mu \|\) 且 \(\| \nu - \nu_0 \| \leqslant \| \mu \|\) ；因此对充分小的 \(\varepsilon\) 有 \(\nu - \mu \in V + V\) ，由此得结论。

推论 3. --- 每个有界正测度 \(\mu\) （\(X\) 上的）都属于由支集有限、含于 \(\mu\) 的支集且范数等于 \(\|\mu\|\) 的正测度所构成的凸集的淡闭包。

与推论 2 相同的推理表明，可以只限于证明 \(\mu\) 属于由支集有限、含于 \(\mathrm{Supp}(\mu)\) 且范数 \(\leqslant \| \mu \|\) 的正测度所构成的凸集 B 的淡闭包。同样，只需证明 \(\mu\) 属于 \(\mathbf{B}^{\circ}\) （即 B 在空间 \(\mathcal{K}(\mathbf{X};\mathbf{R})\) 中的极集）的极集（TVS, II, §6, No.~3, Th. 1）；这就是说，对 \(f\in \mathcal{K}(\mathbf{X};\mathbf{R})\) ，关系 \(\langle f,\varepsilon_a\rangle \geqslant -1 / \| \mu \|\) 对一切 \(a\in \mathrm{Supp}(\mu)\) 成立蕴涵 \(\langle f,\mu \rangle \geqslant -1\) ，而这是 No.~3 命题 8 的推论 2 的推论。

推论 4. --- 在空间 \(\mathcal{M}(\mathrm{X};\mathbf{C})\) 中，点测度所成之集对严格紧收敛拓扑是全的（§1, No.~10）。

在锥 \(\mathcal{M}_{+}(X)\) 上，严格紧收敛拓扑与淡拓扑相同（§1, No.~10, 命题 18），且 X 上的每个测度都可写成 \(\mu_{1}-\mu_{2}+i\mu_{3}-i\mu_{4}\) ，其中诸 \(\mu_{j}\ (1\leqslant j\leqslant4)\) 是正测度；因此结论由 Th. 1 得出。

命题 13. --- 设 \(\mu\) 是局部紧空间 \(X\) 上的一个测度。为使一点 \(x_0\) 属于 \(\operatorname{Supp}(\mu)\) ，必须且只须点测度 \(\varepsilon_{x_0}\) 属于测度 \(g \cdot \mu\) 所成之集的淡闭包，其中 \(g\) 跑遍满足 \(\| g \cdot \mu \| \leqslant 1\) 的具紧支集连续函数所成之集。

依据 No.~2 的命题 6，该条件显然是充分的。为证其必要性，设 \(x_0 \in \operatorname{Supp}(\mu)\) ；考虑有限个函数 \(f_k\) （ \(1 \leqslant k \leqslant n\) ，属于 \(\mathcal{K}(X; \mathbf{C})\) ）及任意一个数 \(\delta > 0\) ；我们要证明存在一个函数 \(g \in \mathcal{K}(X; \mathbf{C})\) ，使 \(\| g \cdot \mu \| \leqslant 1\) 且

\[
\left| f _ {k} \left(x _ {0}\right) - \mu \left(g f _ {k}\right) \right| \leqslant \delta
\]

对 \(1 \leqslant k \leqslant n\) 成立。设 U 是 \(x_0\) 的一个相对紧开邻域，使诸 \(f_k\) （ \(1 \leqslant k \leqslant n\) ）中每一个在 U 上的振幅 \(\leqslant \delta/2\) 。按假设，由于 \(x_0 \in \text{Supp}(\mu)\) ，存在一个函数 \(g_0 \in \mathcal{K}(X; \mathbf{C})\) ，其支集含于 U 且使 \(\mu(g_0) \neq 0\) ；测度 \(\nu = g_0 \cdot \mu\) 不为零，因为对每个在 U 上等于 1 的函数 \(f \in \mathcal{K}(X; \mathbf{C})\) ，有 \(\nu(f) = \mu(g_0) \neq 0\) 。此外， \(\nu\) 是有界的（No.~3, 命题 11）；

用标量乘 \(g_{0}\) ，可设 \(\|\nu\|=1\) 。在此条件下，令 \(\alpha_{k}=f_{k}(x_{0})\) ，则对 \(1\leqslant k\leqslant n\) 及每个函数 \(h\in\mathcal{K}(\mathrm{X};\mathbf{C})\) ，可写

\[
f _ {k} (x _ {0}) - \nu (f _ {k} h) = \alpha_ {k} (1 - \nu (h)) + \nu ((\alpha_ {k} - f _ {k}) h).
\]

由于 \(\nu\) 的支集在 U 中，可把它与其在 U 上的限制等同起来；于是假设 \(\| \nu \| = 1\) 蕴涵存在一个函数 \(h \in \mathcal{K}(\mathrm{X};\mathbf{C})\) ，其支集含于 U，使 \(\| h \| \leqslant 1\) 且 \(|\alpha_{k}(1 - \nu (h))| \leqslant \delta / 2\) 对 \(1 \leqslant k \leqslant n\) 成立。此外，U 的定义表明 \(|(\alpha_{k} - f_{k}(x))h(x)| \leqslant \delta / 2\) 对所有 \(x \in \mathrm{U}\) 成立；由于 \(\| \nu \| = 1\) 且 \(\operatorname{Supp}(\nu) \subset \mathrm{U}\) ，我们有 \(|\nu ((\alpha_{k} - f_{k})h)| \leqslant \delta / 2\) ，于是令 \(g = g_0h\) ，

\[
\left| f _ {k} \left(x _ {0}\right) - \mu \left(g f _ {k}\right) \right| \leqslant \delta \quad \text { for } 1 \leqslant k \leqslant n.
\]

这就证明了该命题，因为 \(\| g\cdot \mu \| = \| (g_0h)\cdot \mu \| \leqslant \| g_0\cdot \mu \| = 1\)

推论. --- 设 \(\mu\) 是 X 上的一个测度。为使 X 上的一个测度 \(\nu\) 属于测度 \(g \cdot \mu\) 所成之集（其中 \(g\) 跑遍 \(\mathcal{K}(X; \mathbf{C})\) ）的淡闭包，必须且只须 \(\operatorname{Supp}(\nu) \subset \operatorname{Supp}(\mu)\) 。

因为依据 No.~3 的命题 10，有 \(\operatorname{Supp}(g\cdot\mu)\subset\operatorname{Supp}(\mu)\) ；因此 \(g\cdot\mu\) 型测度的每个淡极限的支集也含于 \(\operatorname{Supp}(\mu)\) （No.~2, 命题 6）。反之，若 \(\operatorname{Supp}(\nu)\subset\operatorname{Supp}(\mu)\) ，则 \(\nu\) 是支集有限且含于 \(\operatorname{Supp}(\mu)\) 的测度的淡极限（Th. 1），因而由命题 13，它属于测度 \(g\cdot\mu\) 所成之集的淡闭包。

